\section{Coupled period reductions and rationality-dependent sources}
\label{sec:coupled-case}
The constructions in this section illustrate two different ways of
moving beyond a single positive matrix. Coupling produces a small space
of special values, which one can then eliminate by linear algebra.
Alternatively, assuming $\theta$ rational can create endpoint zeros
that were unavailable before the assumption. Both operations are exact;
neither automatically preserves the property needed for an irrationality
proof.

\subsection{An explicit five-period family}\label{subsec:five-period}
\record{P03}
A coupled integral led to the rational series defined below. We use the
series as the definition here, so that reconstructing its coefficients
does not require the original geometric or spherical-integral derivation.
For an integer $n\geq0$, put $\lambda=4n+1$ and
\begin{equation}\label{eq:coupled-A}
 A_j(k)=\frac{(4n)!}{n!}
 \frac{(-1)^j\binom{3n}{j}}
 {\displaystyle\prod_{q=n+j+1}^{4n}(k+q)
  \displaystyle\prod_{q=4n+2}^{4n+j+1}(k+q)},\qquad 0\leq j\leq3n.
\end{equation}
An empty product is one. Define
\[
 L_j=2(k+3n+1+j),\quad M_l=2(k+2n+1+l),\quad
 T_{jl}=2(k+2n+1+j+l),\quad V_{jl}=2(k+1+j+l),
\]
and
\begin{align}
 I_n(k)&=2\sum_{j,l=0}^{3n}A_j(k)A_l(k)
 \left(\frac1{L_jT_{jl}}+\frac1{L_jM_l}+\frac1{M_lV_{jl}}\right),\notag\\
 C_n(k)&=\frac{\prod_{r=1}^{4n}(k+2r)}{2^{4n}(4n)!},\qquad
 R_n(k)=\frac{4(n!)^4C_n(k)I_n(k)}{(4n)!(k+\lambda)^2},\notag\\
 J_n&=\sum_{\ell\geq0}(-1)^\ell R_n(2\ell).
 \label{eq:coupled-R}
\end{align}
Every uncombined summand of $R_n(k)$ is $O_n(k^{-2n-4})$:
each $A_j$ has $3n$ denominator factors, the radial fractions supply
two more, the angular denominator another two, and $C_n$ has degree
$4n$. Thus the series defining $J_n$ converges absolutely. The formulas
are finite rational operations followed by an elementary convergent
sum; they define the family without fitting a recurrence.

\begin{proposition}\label{prop:five-period}
For each $n\geq0$, the formulas above give an explicitly computable
rational decomposition
\begin{equation}\label{eq:five-period}
 J_n=r_n+b_n\beta(2)+a_n\theta+\ell_n\log2+z_n\zeta(3).
\end{equation}
No linear independence of these five numbers is assumed.
\end{proposition}
\begin{proof}
All poles of $R_n$ have the form $-q$, $1\leq q\leq8n+1$, with
order at most four. Away from $q=\lambda$, the two $A$ factors
supply at most two orders and the radial fraction at most two.
At $q=\lambda$ the $A$ factors are regular, so the angular double
pole and the radial fractions still give at most four. Properness
was established above. Consequently
\[
 R_n(k)=\sum_{q=1}^{8n+1}\sum_{s=1}^4\frac{c_{q,s}}{(k+q)^s},
 \qquad c_{q,s}\in\Q.
\]
There is a useful reflection symmetry. Under
$k'=-2\lambda-k$, $j'=3n-j$, formula~\eqref{eq:coupled-A} gives
$A_j(k')=A_{j'}(k)$: the denominator factors reverse, and their signs
cancel the change in $(-1)^j$. The radial factors transform as
$L\leftrightarrow-M$ and $T\leftrightarrow-V$; swapping $j,l$
then preserves their sum. Both $C_n$ and $(k+\lambda)^2$ are
invariant. Hence
\[
 R_n(-2\lambda-k)=R_n(k),\qquad
 c_{2\lambda-q,s}=(-1)^s c_{q,s}.
\]
For $q=2h+1$ and $q=2h+2$, respectively, the required alternating
sums are
\begin{align*}
 \sum_{\ell\geq0}\frac{(-1)^\ell}{(2\ell+2h+1)^s}
 &=(-1)^h\left(\beta(s)-\sum_{j=0}^{h-1}
                          \frac{(-1)^j}{(2j+1)^s}\right),\\
 \sum_{\ell\geq0}\frac{(-1)^\ell}{(2\ell+2h+2)^s}
 &=2^{-s}(-1)^h\left(\eta(s)-\sum_{j=1}^{h}
                          \frac{(-1)^{j-1}}{j^s}\right),
\end{align*}
where $\eta(1)=\log2$ and $\eta(3)=3\zeta(3)/4$.
Across reflected odd poles the tail signs agree, canceling odd beta
orders; across reflected even poles they are opposite, canceling even
eta orders. The central pole is odd and has only even orders.
Only $\beta(2),\beta(4),\eta(1),\eta(3)$ and rational tails remain.
The individual order-one sums converge conditionally; there are only
finitely many of them, so taking their limits is legitimate.
\end{proof}
This proof is also an exact algorithm. At a pole of maximal order four,
expand $(k+q)^4R_n(k)$ through degree three in $k+q$ to recover the
four coefficients, then use the two finite-tail formulas. For example,
\[
 J_0=6\theta,\qquad
 J_1=\frac{301}{384}-\frac{503}{64}\beta(2)
        +\frac{243}{32}\theta-\frac{73}{24}\log2+\frac{27}{32}\zeta(3).
\]
The second displayed coefficient vector has primitive integer form
$(301,-3018,2916,-1168,324)$ in the order
$(1,\beta(2),\theta,\log2,\zeta(3))$.

\paragraph{Elimination and what it does not preserve.}
Choose four indices $n_0,\ldots,n_3$ and form the $3\times4$ rational
matrix whose columns are $(b_{n_j},\ell_{n_j},z_{n_j})^T$.
Let $w_j$ be its signed maximal minors, with signs chosen so that
$\sum_jw_j(b_{n_j},\ell_{n_j},z_{n_j})=0$. Then
\begin{equation}\label{eq:coupled-cofactor}
 \sum_{j=0}^3w_jJ_{n_j}
 =\left(\sum_jw_ja_{n_j}\right)\theta+\sum_jw_jr_{n_j}.
\end{equation}
If the matrix has rank less than three, use an explicitly computed
rational null vector instead; the maximal minors would all vanish.
Thus the auxiliary constants can be eliminated without assuming they
are independent. But positive individual integrals can cancel in this
signed sum, and the minors can have large denominators. To use
\eqref{eq:coupled-cofactor} one still needs nonvanishing and a bound
after dividing out the full common content of its two coefficients.
The four recorded finite cofactor experiments do not supply those
all-index estimates. The fixed period-space theorem is the reusable
conclusion here; it is not a proof that the final linear form is small.

\subsection{Endpoint zeros obtained by assuming rationality}
\label{subsec:niven-case}\record{E07}
Niven-type proofs suggest a different starting point: use a hypothetical
rational value to construct a positive function with many endpoint zeros.
Let
\[
 \mathcal F(x)=\sum_{m\geq1}\frac{\chi_4(m)x^m}{m^4},\qquad
 E=x\frac{d}{dx},\qquad E^4\mathcal F=\frac{x}{1+x^2}.
\]
The integral representation
\[
 \mathcal F(x)=\frac{x}{6}\int_0^1
                       \frac{(-\log t)^3}{1+x^2t^2}\dd t
\]
shows that $\mathcal F'(x)>0$ on $[0,1]$: its differentiated
integrand contains $(1-x^2t^2)/(1+x^2t^2)^2$, positive almost everywhere.
Assume $\theta=u/v$ with positive integers $u,v$, and define
\begin{equation}\label{eq:niven-source}
 A_n(x)=\frac{[v\mathcal F(x)(u-v\mathcal F(x))]^n}{n!},\qquad
 I_n=\int_0^1\frac{A_n(x)}{1+x^2}\dd x.
\end{equation}
Then $0<I_n\leq(u^2/4)^n/n!\to0$. The source has a zero of order
$n$ at each endpoint. An integrality theorem for an unbounded sequence
of these $I_n$ would therefore give a contradiction. The issue is that
the endpoint zeros do not provide the requisite integrality theorem.

\begin{proposition}\label{prop:niven-prime}
For any fixed $n\geq1$ and nonzero integers $u,v$, every odd prime
$p>n$ with $p\nmid uv$ satisfies
\[
 v_p\bigl(A_n^{(p+n-1)}(0)\bigr)=-3.
\]
Under the rationality assumption the first nonzero Euler derivative
at the other endpoint is
\[
 E^nA_n(1)=(-uv)^n\left(\frac{\pi^3}{32}\right)^n.
\]
\end{proposition}
\begin{proof}
Expand
\[
 A_n=\frac1{n!}\sum_{k=0}^n(-1)^k\binom nk
                   u^{n-k}v^{n+k}\mathcal F^{n+k}.
\]
At degree $N=p+n-1<2p$, a $p$-denominator can only come from the
coefficient $\chi_4(p)/p^4$ in one factor. In $\mathcal F^{n+k}$,
using that coefficient requires degree at least $p+n+k-1$.
Only $k=0$ contributes, with $n$ choices for that factor and the
linear coefficient in every other factor. Its contribution is
$u^nv^n\chi_4(p)/((n-1)!p^4)$, of valuation $-4$; all other terms
are $p$-integral. Multiplying by $N!$, whose valuation is one,
proves the claim. At $x=1$, use the coordinate $t=\log x$.
The linear term of $v\mathcal F(e^t)(u-v\mathcal F(e^t))$ is
$-uv(E\mathcal F)(1)t$. Since $(E\mathcal F)(1)=\beta(3)=\pi^3/32$,
the coefficient of $t^n$ gives the second formula.
\end{proof}
Thus not all origin derivatives can be integral, and rationality of
the endpoint does not rationalize its derivative. This excludes a
blanket endpoint-integrality argument, while leaving open constructions
that use only selected derivatives and control their denominators.

\subsection{Changing variables and the differential obstruction}
In the coordinate $z=\mathcal F(x)$, the source becomes a polynomial,
but the natural corresponding integral is explicitly
\begin{equation}\label{eq:niven-beta}
 \int_0^1 A_n(x)\mathcal F'(x)\dd x
       =\frac{u^{2n+1}n!}{v(2n+1)!}.
\end{equation}
Indeed substitute $w=v\mathcal F(x)$ and use
$\int_0^u w^n(u-w)^n\dd w=u^{2n+1}(n!)^2/(2n+1)!$.
This is a positive rational number. For
$n=(p-1)/2$, $p\nmid uv$, its valuation at the odd prime $p$ is
$-1$. Clearing its denominator and removing common numerator content
leaves the constant polynomial $1$, rather than shrinking integers.

One cannot instead keep the original Euler derivative and treat
$\mathcal F$ as its only additional rational-function coordinate.
In fact
\begin{equation}\label{eq:niven-field}
 E\mathcal F\notin\C(x,\mathcal F).
\end{equation}
Here is a direct functional proof. The third Euler derivative is
$\arctan x$; continuation once around $x=i$ changes it by $\pi$.
Integrating with $E$, and taking the lower changes to vanish at $i$,
gives
\[
 \Delta\mathcal F=\frac\pi6\log^3(x/i),\qquad
 \Delta(E\mathcal F)=\frac\pi2\log^2(x/i).
\]
If $E\mathcal F=R(x,\mathcal F)$ with $R\in\C(x,z)$, repeated
continuation at a generic nearby $x$ would give
$R(x,\mathcal F+m\Delta\mathcal F)
 =E\mathcal F+m\Delta(E\mathcal F)$ for every integer $m$.
A rational function of $z$ agreeing with an affine function on this
infinite set must itself be affine. Its slope would be
$3/\log(x/i)$, which is not rational in $x$. This proves
\eqref{eq:niven-field}. The larger field containing all four Euler
primitives does close under differentiation. It would require a new
endpoint-elimination and arithmetic argument, which this calculation
does not exclude.

\paragraph{A useful control example.}
The failure is not repaired by imposing rational Taylor coefficients,
positive rational forcing and a nonpolynomial primitive. The function
\[
 \mathcal H(x)=x+\frac{\Li_4(x/2)-\Li_4(x^2/2)}{32}
\]
has $\mathcal H(0)=0$, $\mathcal H(1)=1$, and
\[
 \frac{E^4\mathcal H}{x}
    =1+\frac1{32}\left(\frac1{2-x}-\frac{16x}{2-x^2}\right)
       \geq\frac{33}{64},\qquad
 \mathcal H'(x)\geq\frac{251}{256}\quad(0\leq x\leq1).
\]
For the derivative estimate use
$z\leq\Li_3(z)\leq z+z^2/[8(1-z)]$ and differentiate the displayed
formula; the two polylogarithmic terms divided by $x$ are bounded
below by $1/2$ and above by $9/8$, respectively. All Taylor
coefficients are rational and the forcing has no pole on the interval.
Nevertheless its endpoint is the rational number $1$.
Thus those structural properties alone cannot supply the missing
irrationality implication. The target's particular pole geometry
would have to enter an additional argument.
